\chapter{HIV/AIDS}
\index{HIV/AIDS}

\section*{Definition and Overview}
HIV (human immunodeficiency virus) is a virus that attacks the immune system, specifically the CD4 white blood cells that coordinate the body's defences. Without treatment, the virus progressively destroys these cells, leaving the body vulnerable to infections and cancers. AIDS (acquired immunodeficiency syndrome) is the most advanced stage of HIV infection, defined by a very low CD4 count or the occurrence of particular opportunistic infections and cancers. HIV is transmitted through blood, sexual fluids, and from mother to child. There is no cure, but modern antiretroviral treatment controls the virus completely, allowing people with HIV to live long, healthy lives and preventing transmission.

\section*{Epidemiology}
Around 30,000 people live with HIV, and around 90\% are diagnosed and receiving treatment. New diagnoses have fallen substantially, particularly among men who have sex with men, reflecting the success of testing, treatment, and prevention including PrEP. Globally, around 39 million people live with HIV, and millions of lives have been saved by antiretroviral therapy. In many countries, mother-to-child transmission is now extremely rare, and people on effective treatment cannot transmit the virus to sexual partners ("undetectable equals untransmittable").

\section*{Aetiology and Pathophysiology}
\begin{itemize}
    \item \textbf{Transmission:} unprotected sex, sharing needles, mother-to-child during pregnancy, birth, or breastfeeding, and contaminated blood or blood products
    \item \textbf{Pathophysiology:} HIV infects CD4 T-lymphocytes, uses them to replicate, and destroys them
    \item \textbf{Acute phase:} rapid viral replication and a flu-like illness weeks after infection
    \item \textbf{Chronic phase:} the virus replicates at lower levels for years while CD4 counts gradually fall
    \item \textbf{AIDS:} when CD4 counts drop below 200 cells/mm³ or opportunistic infections occur
    \item \textbf{Opportunistic infections:} pneumocystis pneumonia, tuberculosis, and others exploit immune failure
    \item \textbf{Antiretroviral therapy:} suppresses viral replication, allowing immune recovery
\end{itemize}

\section*{Clinical Features}
\begin{itemize}
    \item Acute seroconversion illness: fever, sore throat, rash, swollen glands, and fatigue, 2--6 weeks after exposure
    \item A long asymptomatic phase lasting years
    \item Persistent swollen lymph nodes
    \item Non-specific symptoms: fatigue, weight loss, night sweats, and recurrent infections
    \item AIDS-defining illnesses: pneumocystis pneumonia, candidiasis of the oesophagus, tuberculosis, and Kaposi sarcoma
    \item Neurological problems: memory loss and peripheral neuropathy
    \item Many people are diagnosed before any symptoms through testing
\end{itemize}

\section*{Differential Diagnosis}
\begin{itemize}
    \item Glandular fever and other viral infections mimic seroconversion illness
    \item Chronic fatigue and weight loss have many causes
    \item Immunodeficiency from other causes
    \item Tuberculosis and other infections occur at all stages
    \item HIV testing is definitive and is recommended routinely in at-risk groups and in pregnancy
\end{itemize}

\section*{When to Seek Medical Care}
Anyone who may have been exposed to HIV should be tested, including after a condom break, needlestick, or sharing of injecting equipment. Post-exposure prophylaxis (PEP) must be started within 72 hours of exposure.

\textbf{Contact your local emergency services for:}
\begin{itemize}
    \item Difficulty breathing, confusion, or seizures in someone with HIV
    \item High fever with a severe headache or stiff neck
    \item Vomiting blood, severe abdominal pain, or collapse
    \item Thoughts of self-harm or suicide
    \item Any emergency in an immunocompromised person
\end{itemize}

\section*{Diagnosis and Investigations}
\begin{itemize}
    \item \textbf{HIV antibody/antigen test:} the standard screening test; most tests detect infection within 4 weeks
    \item \textbf{Rapid tests:} finger-prick testing available in clinics and at home
    \item \textbf{Confirmatory testing:} positive screens are confirmed with additional tests
    \item \textbf{Viral load:} measures the amount of virus and guides treatment
    \item \textbf{CD4 count:} assesses immune status
    \item \textbf{Drug resistance testing:} before starting treatment
    \item \textbf{Screening for other infections:} hepatitis, syphilis, and tuberculosis
\end{itemize}

\section*{Management}
\begin{itemize}
    \item \textbf{Antiretroviral therapy:} a combination of medicines taken daily, recommended for everyone with HIV from diagnosis
    \item \textbf{Treatment goals:} undetectable viral load and immune recovery
    \item \textbf{Undetectable = untransmittable:} people with an undetectable viral load cannot pass HIV to sexual partners
    \item \textbf{Monitoring:} regular viral load, CD4 count, and side-effect checks
    \item \textbf{Opportunistic infection prevention:} vaccination, preventive medicines, and prompt treatment
    \item \textbf{Treatment in pregnancy:} prevents transmission to the baby, which is now rare in many countries
    \item \textbf{Support:} HIV services, peer support, and mental health care
    \item \textbf{Healthy living:} smoking cessation, diet, and cardiovascular risk management
\end{itemize}

\section*{Complications}
\begin{itemize}
    \item Opportunistic infections and AIDS-defining cancers in untreated disease
    \item Chronic inflammation-related conditions: cardiovascular disease, kidney disease, and osteoporosis
    \item Neurocognitive impairment
    \item Medication side effects and interactions
    \item Stigma, discrimination, and mental health difficulties
    \item Transmission to partners or children if untreated
\end{itemize}

\section*{Prognosis and Outcomes}
With modern treatment, the outlook for HIV is excellent. People diagnosed early and treated consistently have a life expectancy close to that of the general population, and an undetectable viral load means they cannot transmit the virus. The main challenges are adherence, side effects, and managing stigma. Without treatment, HIV progresses to AIDS and is fatal within years, which is why testing and early treatment are so important. Prevention, including PrEP and condoms, is highly effective.

\section*{Prevention and Self-Care}
\begin{itemize}
    \item Use condoms consistently
    \item Consider PrEP (pre-exposure prophylaxis) if at risk of HIV
    \item Use PEP within 72 hours of a potential exposure
    \item Never share needles or injecting equipment
    \item Get tested regularly if at risk
    \item Take antiretroviral treatment exactly as prescribed
    \item Pregnant women are tested so transmission can be prevented
    \item Seek support for stigma, disclosure, and mental health
\end{itemize}

\section*{Sources and Further Reading}
The following reputable sources provide further information for healthcare students and patients seeking reliable, current advice about HIV/AIDS.

\begin{itemize}
    \item Healthdirect. \textit{HIV and AIDS: overview.} https://www.healthdirect.gov.au/hiv-aids
    \item Positive Life NSW and other PLHIV organisations. \textit{Support and information.} https://www.positivelife.org.au/
    \item ASHM. \textit{HIV testing and treatment guidelines.} https://www.ashm.org.au/
    \item World Health Organization. \textit{HIV/AIDS: fact sheet.} https://www.who.int/
    \item UNAIDS. \textit{Global HIV data.} https://www.unaids.org/
    \item NHS. \textit{HIV and AIDS: overview.} https://www.nhs.uk/conditions/hiv-and-aids/
\end{itemize}
